\subsection{RQ1: End-to-End Effectiveness}
\label{sec:rq1}

% Notes: docs/exhibits.md#tab_rq1_end_to_end_summary
\begin{table}[t]
  \centering
  \setlength{\abovecaptionskip}{2pt}%
  \setlength{\belowcaptionskip}{3pt}%
  \caption{End-to-end effectiveness of our controller (RQ1).}
  \label{tab:rq1_preventive_alignment}
  \label{tab:rq1_controller_behavior}
  {\scriptsize
  \setlength{\arrayrulewidth}{0.35pt}%
  \setlength{\doublerulesep}{0.8pt}%
  \renewcommand{\arraystretch}{1.04}%
  \newlength{\rqsurvivedwidth}%
  \settowidth{\rqsurvivedwidth}{\textbf{10/10} (--)}%
  \newcommand{\rqsurvived}[1]{\makebox[\rqsurvivedwidth][r]{#1}}%
  \newlength{\rqstageheadingwidth}%
  \settowidth{\rqstageheadingwidth}{\textbf{Stage execution}}%
  \newsavebox{\rqtabbox}%
  \newcommand{\rqtabular}{%
  \begin{tabular}{|c|c|c|r|r|r|r|r|r|}
    \hline\hline
    \multirow[c]{3}{*}[-0.5\normalbaselineskip]{\textbf{Device}} &
    \multirow[c]{3}{*}[-0.5\normalbaselineskip]{\textbf{Condition}} &
    \multirow[c]{3}{*}[-0.5\normalbaselineskip]{\makecell[c]{\textbf{Starting}\\\textbf{overheat}\\\textbf{level}}} &
    \multicolumn{2}{c|}{\textbf{Deadline safety}} &
    \multicolumn{4}{c|}{\textbf{Our controller}} \\
    \cline{4-9}
    & & &
    \multicolumn{2}{c|}{\textbf{Completed runs}} &
    \multicolumn{2}{c|}{\textbf{Stage execution}} &
    \multicolumn{2}{c|}{\textbf{Pacing cost}} \\
    \cline{4-9}
    & & &
    \multicolumn{1}{c|}{\textbf{Baseline}} &
    \multicolumn{1}{c|}{\textbf{Ours}} &
    \multicolumn{1}{c|}{\makecell[c]{\makebox[.5\rqstageheadingwidth]{\boldmath$M$}\\\textbf{(\%)}}} &
    \multicolumn{1}{c|}{\makecell[c]{\makebox[.5\rqstageheadingwidth]{\boldmath$S$}\\\textbf{(\%)}}} &
    \multicolumn{1}{c|}{\makecell[c]{\textbf{Activation}\\\textbf{(\%)}}} &
    \multicolumn{1}{c|}{\makecell[c]{\textbf{Delay}\\\textbf{(ms)}}} \\
    \hline\hline
    \multirow{14}{*}{\makecell[c]{S26\\Ultra}}
      & \multirow{7}{*}{\makecell[c]{12MP\\normal}}
      & Lv0 & \rqsurvived{\textbf{10/10} (--)} & \rqsurvived{\textbf{10/10} (--)}
        & 100.0 & 100.0 & 0.0 & -- \\
    \cline{3-9}
      & & Lv1 & \rqsurvived{\textbf{10/10} (--)} & \rqsurvived{\textbf{10/10} (--)}
        & 98.0 & 100.0 & 0.7 & 127 \\
    \cline{3-9}
      & & Lv2 & \rqsurvived{7/10 (24)} & \rqsurvived{\textbf{10/10} (--)}
        & 90.3 & 100.0 & 12.4 & 262 \\
    \cline{3-9}
      & & Lv3 & \rqsurvived{1/10 (13)} & \rqsurvived{\textbf{10/10} (--)}
        & 85.3 & 98.0 & 35.9 & 284 \\
    \cline{3-9}
      & & Lv4 & \rqsurvived{0/10 (8)} & \rqsurvived{\textbf{10/10} (--)}
        & 58.0 & 95.7 & 37.9 & 437 \\
    \cline{3-9}
      & & Lv5 & \rqsurvived{0/10 (8)} & \rqsurvived{\textbf{10/10} (--)}
        & 45.3 & 94.3 & 42.1 & 470 \\
    \cline{3-9}
      & & Lv6 & \rqsurvived{0/10 (7)} & \rqsurvived{\textbf{10/10} (--)}
        & 36.3 & 88.3 & 34.1 & 521 \\
    \Xcline{2-9}{0.7pt}
      & \multirow{7}{*}{\makecell[c]{24MP\\CPU\\stress}}
      & Lv0 & \rqsurvived{7/10 (26)} & \rqsurvived{\textbf{10/10} (--)}
        & 92.7 & 100.0 & 37.9 & 338 \\
    \cline{3-9}
      & & Lv1 & \rqsurvived{2/10 (19)} & \rqsurvived{\textbf{10/10} (--)}
        & 66.7 & 90.0 & 35.2 & 470 \\
    \cline{3-9}
      & & Lv2 & \rqsurvived{0/10 (14)} & \rqsurvived{\textbf{10/10} (--)}
        & 58.3 & 79.0 & 37.2 & 377 \\
    \cline{3-9}
      & & Lv3 & \rqsurvived{0/10 (5)} & \rqsurvived{\textbf{10/10} (--)}
        & 35.0 & 54.3 & 37.9 & 418 \\
    \cline{3-9}
      & & Lv4 & \rqsurvived{0/10 (3)} & \rqsurvived{\textbf{10/10} (--)}
        & 44.7 & 77.0 & 36.6 & 400 \\
    \cline{3-9}
      & & Lv5 & \rqsurvived{0/10 (4)} & \rqsurvived{\textbf{10/10} (--)}
        & 17.3 & 52.7 & 29.3 & 367 \\
    \cline{3-9}
      & & Lv6 & \rqsurvived{0/10 (3)} & \rqsurvived{\textbf{10/10} (--)}
        & 25.0 & 46.7 & 29.3 & 525 \\
    \Xcline{1-9}{0.7pt}
    \multirow{14}{*}{S26}
      & \multirow{7}{*}{\makecell[c]{12MP\\normal}}
      & Lv0 & \rqsurvived{7/10 (28)} & \rqsurvived{\textbf{10/10} (--)}
        & 100.0 & 100.0 & 9.0 & 159 \\
    \cline{3-9}
      & & Lv1 & \rqsurvived{5/10 (26)} & \rqsurvived{\textbf{10/10} (--)}
        & 100.0 & 100.0 & 12.8 & 241 \\
    \cline{3-9}
      & & Lv2 & \rqsurvived{4/10 (22)} & \rqsurvived{\textbf{10/10} (--)}
        & 100.0 & 100.0 & 18.3 & 256 \\
    \cline{3-9}
      & & Lv3 & \rqsurvived{0/10 (15)} & \rqsurvived{\textbf{10/10} (--)}
        & 89.7 & 100.0 & 40.7 & 277 \\
    \cline{3-9}
      & & Lv4 & \rqsurvived{0/10 (9)} & \rqsurvived{\textbf{10/10} (--)}
        & 56.0 & 74.7 & 53.1 & 297 \\
    \cline{3-9}
      & & Lv5 & \rqsurvived{0/10 (9)} & \rqsurvived{\textbf{10/10} (--)}
        & 45.7 & 79.0 & 43.8 & 392 \\
    \cline{3-9}
      & & Lv6 & \rqsurvived{0/10 (9)} & \rqsurvived{\textbf{10/10} (--)}
        & 41.3 & 66.3 & 52.1 & 410 \\
    \Xcline{2-9}{0.7pt}
      & \multirow{7}{*}{\makecell[c]{24MP\\CPU\\stress}}
      & Lv0 & \rqsurvived{1/10 (23)} & \rqsurvived{\textbf{10/10} (--)}
        & 100.0 & 100.0 & 52.4 & 330 \\
    \cline{3-9}
      & & Lv1 & \rqsurvived{1/10 (21)} & \rqsurvived{\textbf{10/10} (--)}
        & 93.3 & 95.3 & 42.8 & 255 \\
    \cline{3-9}
      & & Lv2 & \rqsurvived{0/10 (19)} & \rqsurvived{\textbf{10/10} (--)}
        & 59.3 & 76.7 & 49.0 & 419 \\
    \cline{3-9}
      & & Lv3 & \rqsurvived{0/10 (9)} & \rqsurvived{\textbf{10/10} (--)}
        & 71.7 & 97.0 & 60.0 & 466 \\
    \cline{3-9}
      & & Lv4 & \rqsurvived{0/10 (6)} & \rqsurvived{\textbf{10/10} (--)}
        & 74.0 & 91.3 & 61.0 & 458 \\
    \cline{3-9}
      & & Lv5 & \rqsurvived{0/10 (5)} & \rqsurvived{\textbf{10/10} (--)}
        & 24.0 & 30.7 & 25.2 & 538 \\
    \cline{3-9}
      & & Lv6 & \rqsurvived{0/10 (6)} & \rqsurvived{\textbf{10/10} (--)}
        & 21.7 & 26.0 & 22.1 & 467 \\
    \hline\hline
  \end{tabular}}%
  \fittabcolsep{\rqtabbox}{\rqtabular}{\columnwidth}{18}%
  \rqtabular\par
  \raggedright
  \textbf{Completed runs} -- How many of the ten runs finished all 30 captures.\newline
  The number in parentheses identifies the earliest capture number at which any run timed out.\par
  \textbf{Stage execution} -- For each stage (\(M\) or \(S\)), the percentage of captures in which it ran.\par
  \textbf{Pacing cost} -- Activation is the percentage of transitions paced. Delay is their median delay.\par
  }
\end{table}


Table~\ref{tab:rq1_controller_behavior} shows that our controller completed every run without Capture Timeout in all 28 settings while still executing optional stages, at the cost of pacing delays.
Each row reports one setting, with stage-execution and pacing-activation percentages computed over all 30 captures of its ten analyzed runs.

\myparagraph{Deadline Safety}
Our controller completed all 280 runs without Capture Timeout, whereas Baseline completed only 55 of 280 runs.
In particular, Baseline timed out in 14 of the 16 settings starting below the thermal guard's level-4 cutoff.
In contrast, our controller avoided these timeouts, and watchdog fallback occurred during \(M\) in only 2 of the 280 runs (0.7\%), both under \textcolor{red}{CPU stress}.

\myparagraph{Stage Execution}
Over all 30 captures, stage execution generally decreased as the starting overheat level increased for each device and capture condition.
Notably, in the twelve settings starting at or above the guard's level-4 cutoff, \(M\) still executed on 17.3--74.0\% of captures.
To examine camera usability during early consecutive captures, we further analyzed the first five captures of each run.
In these captures, the average execution counts per run were at least 4.5 for \(M\) and 4.6 for \(S\) in every setting.
These results indicate that our controller executed optional stages even in settings where the static thermal guard would skip them at the first capture.

\myparagraph{Pacing Cost}
On average, pacing activated on 33.9\% of the 29 shot-to-shot transitions, with median delays of 127--538~ms in settings where pacing activated.
Even with these delays, the mean shot-to-shot interval was 21--54\% of the estimated interval with parallel capture disabled.
This estimate is the time from request acceptance to the completion of a draft sequence that executes both \(M\) and \(S\) in the same setting.
Moreover, during the first five captures, pacing did not activate in any of the eight production-default 12MP normal settings starting below the guard's level-4 cutoff.
These findings show that our controller retained the responsiveness benefit of parallel capture despite pacing delays.